\documentclass[12pt]{article}

\usepackage[margin=2.5cm]{geometry}
\usepackage{listings}
\usepackage{amsmath}
\usepackage{amssymb}
\usepackage{courier}
\usepackage{color}

\definecolor{mygreen}{rgb}{0, 0.6, 0}
\definecolor{mygray}{rgb}{0.5, 0.5, 0.5}
\definecolor{mymauve}{rgb}{0.58, 0, 0.82}
\definecolor{stringliterals}{RGB}{245, 141, 66}
\definecolor{background}{RGB}{255, 249, 236}
\definecolor{keywords}{RGB}{2, 166, 199}
\definecolor{comments}{RGB}{150, 0, 0}

\lstset{ %
  backgroundcolor=\color{background},   % choose the background color; you must add \usepackage{color} or \usepackage{xcolor}
  basicstyle=\footnotesize\ttfamily,      % the size of the fonts that are used for the code
  breakatwhitespace=false,         % sets if automatic breaks should only happen at whitespace
  breaklines=true,                 % sets automatic line breaking
  captionpos=b,                    % sets the caption-position to bottom
  commentstyle=\color{mygray},    % comment style
  deletekeywords={...},            % if you want to delete keywords from the given language
  escapeinside={\%*}{*)},          % if you want to add LaTeX within your code
  extendedchars=true,              % lets you use non-ASCII characters; for 8-bits encodings only, does not work with UTF-8
  frame=single,                    % adds a frame around the code
  keepspaces=true,                 % keeps spaces in text, useful for keeping indentation of code (possibly needs columns=flexible)
  keywordstyle=\color{keywords},       % keyword style
  language=Python,                 % the language of the code
  morekeywords={None, self, as},               % if you want to add more keywords to the set
  numbers=none,                    % where to put the line-numbers; possible values are (none, left, right)
  numbersep=0pt,                   % how far the line-numbers are from the code
  numberstyle=\tiny\color{mygray}, % the style that is used for the line-numbers
  rulecolor=\color{black},         % if not set, the frame-color may be changed on line-breaks within not-black text (e.g. comments (green here))
  showspaces=false,                % show spaces everywhere adding particular underscores; it overrides 'showstringspaces'
  showstringspaces=false,          % underline spaces within strings only
  showtabs=false,                  % show tabs within strings adding particular underscores
  stepnumber=2,                    % the step between two line-numbers. If it's 1, each line will be numbered
  stringstyle=\color{stringliterals},     % string literal style
  tabsize=4,                       % sets default tabsize to 2 spaces
  title=\lstname                   % show the filename of files included with \lstinputlisting; also try caption instead of title
}

\title{Entrega - Modelos Estocásticos}
\author{Francisco Peláez Arias}
\date{\today}

\begin{document}

\maketitle

\begin{lstlisting}[title=src/simulation.py]
    from dataclasses import dataclass
    import numpy as np
    from matplotlib import pyplot as plt  

    @dataclass
    class SimulationResult:
        paths: list[list[int]]
        estimated_final_distribution: np.ndarray

    class Simulation:
        def __init__(
            self,
            trials: int,
            steps: int,
            n_states: int,
            initial_distribution: np.ndarray,
            transition_matrix: np.ndarray,
        ) -> None:
            self.trials = trials
            self.steps = steps
            self.state_space = [i for i in range(n_states)]
            self.n_states = n_states
            self.initial_distribution = initial_distribution
            self.transition_matrix = transition_matrix
            self.result: None | SimulationResult = None

\end{lstlisting}

\begin{lstlisting}[title={src/simulation.py, después de \_\_init\_\_()}]
        def run(self) -> None:
            ...
            paths = self._compute_paths()
            est_final_dist = self._estimate_final_distribution()
            self.result = SimulationResult(paths, est_final_dist)
            ...
\end{lstlisting}

\vspace{1.3cm}

\begin{lstlisting}[title={src/simulation.py, después de run()}]
        def _compute_paths(self) -> list[list[int]]:
            paths = [[] for _ in range(self.trials)]
            for i in range(self.trials):
                path = [0 for _ in range(self.steps)]
                in_state = np.random.choice(
                    self.state_space,
                    p=self.initial_distribution
                )
                path[0] = in_state
                pv_state = in_state
                for j in range(1, self.steps):
                    cr_state = np.random.choice(
                        self.state_space,
                        p=self.transition_matrix[pv_state]
                    )
                    path[j] = cr_state
                    pv_state = cr_state
                paths[i] = path
          return paths
\end{lstlisting}

\begin{lstlisting}[title={src/simulation.py, después de \_compute\_paths()}]
        def _estimate_final_distribution(
            self, trials: int = 1_000
        ) -> np.ndarray:
            return np.concatenate(
                [
                    self.estimate_distribution_from_initial_state(
                        i, self.steps, trials
                    )
                    for i in range(trials)
                ]
            )
\end{lstlisting}

\begin{lstlisting}[title={src/simulation.py, después de \_estimate\_final\_distribution()}]
        def estimate_final_distribution_from_initial_state(
            self,
            initial_state: int,
            final_epoch: int,
            trials: int = 1_000,
        ) -> np.ndarray:
            counts = np.zeros((self.n_states,))
            for _ in range(trials):
                state = initial_state
                for _ in range(1, final_epoch):
                    state = np.random.choice(
                        self.state_space, p=self.transition_matrix[state]
                    )
                counts[state] += 1
            return counts / counts.sum()
\end{lstlisting}

\end{document}